\subsection{Experimental Setup}
\label{sec:ExperimentalSetup}
All our experiments were run on Intel Xeon Gold 6248R processors running at
3.00 GHz with an NVIDIA A100 80 GiB GPU.
We use the Qwen3 4B as the base model \citep{qwen3-tr}. Specifically, we use the
Qwen3 4B Thinking 2507 version, which was better at following formatting
instructions during our early experiments. We use temperature $0.6$ and a
maximum context length of 32\,768 tokens during inference.
To fine-tune our models, we use low-rank adaptation (LoRA)
\citep{hu-et-al-iclr2022}. The hyperparameters used for fine-tuning can be found
in Appendix~\ref{sec:hyperparameters}. To avoid stacking up or retraining existing adapters,
when we train generation $n$, we fine-tune the base model Qwen3 4B with the
traces from generations $0, 1, 2, \dots, n-1$. To validate the computed plans, we use
VAL \citep{howey-long-icaps2003wscompetition}, which is commonly used in
planning competitions.


\paragraph{Prompt.}

Our prompt format is based on the prompt used by \citet{correa-et-al-neurips2025}
for end-to-end PDDL planning. To solve a task $T$ from a domain $D$, the
prompt contains the following components:

\begin{enumerate}
\item The following instruction: ``\textit{You are generating plans for PDDL tasks. You
  will be given the PDDL domain and the PDDL instance, and you need to return
  the plan.}''
\item Two examples of PDDL domains, tasks, and their respective plans.
\item The PDDL domain $D$ and the PDDL task $T$.
\end{enumerate}
We use tasks from the traditional Gripper and Logistics domains as
examples. These domains are not included in our experiments. The plans used
to illustrate the examples are not optimal plans.



\subsection{Fine-Tuning Hyperparameters}
\label{sec:hyperparameters}

Below we include the hyperparameters used during our supervised fine-tuning with
LoRA. At each generation, we randomly select 10$\%$ of the traces to be used as
validation set. The validation loss is evaluation every 10 steps.

\begin{center}
\begin{tabular}{lr}
\toprule
\textbf{Hyperparameter}      & \textbf{Value}     \\ \midrule
 Maximum Context Length      & 32\,768            \\
 Batch Size                  & 1                  \\
 Gradient Accumulation Steps & 1                  \\
 Epochs                      & 2                  \\
 Optimizer                   & AdamW (fused)      \\
 Learning Rate               & $1 \times 10^{-5}$ \\
 Learning Rate Scheduler     & Cosine             \\
 Warm-up Ratio                & 0.02               \\
 Weight Decay                & 0.01               \\
 Max. Gradient Norm          & 1.0                \\
 LoRA Rank                   & 16                 \\
 LoRA Alpha                  & 32                 \\
 LoRA Dropout                & 0.05               \\
 LoRA Bias                    & None               \\
\bottomrule
\end{tabular}
\end{center}
